\documentclass[11pt]{article}
\usepackage{preamble}

\newcommand{\lessonrow}[2]{\textbf{#1} & #2 \\ \hline}

\begin{document}

\packetheading{Lesson 3-5 \textperiodcentered\ Period 2 \textperiodcentered\ Student Packet}{Multiplicity and Factored-Form Graphing}

\sectionbanner{OBJECTIVES}
{\small
\renewcommand{\arraystretch}{1.25}
\noindent\begin{tabular}{|p{1.8in}|p{4.8in}|}
\hline
\lessonrow{Math Objective}{Use the multiplicity of each zero to determine whether the graph of a polynomial crosses the x-axis or is tangent to (turns back from) it at each zero.}
\lessonrow{Language Objective}{Students will use the frame ``The zero x = \_\_\_ has multiplicity \_\_\_, so the graph \_\_\_ the x-axis.'' to describe each zero's graph behavior.}
\lessonrow{Essential Understanding}{A factored-form polynomial reveals its zeros directly. The multiplicity of each zero --- odd or even --- controls whether the graph crosses or is tangent to the x-axis at that zero.}
\end{tabular}
}

\sectionbanner{TOPIC GOALS / ESSENTIAL QUESTION / MATERIALS}
{\small
\renewcommand{\arraystretch}{1.25}
\noindent\begin{tabular}{|p{1.8in}|p{4.8in}|}
\hline
\lessonrow{Topic Goals}{Students apply the multiplicity rule to four factored-form polynomials, one per practice item. Every item is drawn from Savvas Lesson 3-5.}
\lessonrow{Essential Question}{How does the multiplicity of each zero determine the graph's behavior at that x-intercept?}
\lessonrow{Materials}{Pencils; student packet; calculator (optional). No Chromebooks required for Period 2.}
\end{tabular}
}

\sectionbanner{LAUNCH --- Multiplicity Rule (teacher models on screen)}

\begin{callout}{\IconBook\ MULTIPLICITY RULES (keep printed on your page)}{calloutgreen}
\textbullet\ If a factor appears ONCE or an ODD number of times, the zero has ODD multiplicity \(\rightarrow\) the graph CROSSES the x-axis at that zero.\\
\textbullet\ If a factor appears an EVEN number of times, the zero has EVEN multiplicity \(\rightarrow\) the graph is TANGENT to the x-axis (touches and turns back) at that zero.\\
\textbullet\ Example: \(f(x) = (x - 2)^3(x + 1)^2\) \(\rightarrow\) \(x = 2\) has multiplicity 3 (crosses); \(x = -1\) has multiplicity 2 (tangent).
\end{callout}

\sectionbanner{EXPLORE --- Think-Pair-Share (35+ min)}
\textit{Work with a partner. Use the multiplicity rules above for every problem.}

\textbf{Bank-item labels (in order):}
\begin{enumerate}[leftmargin=1.6em,itemsep=2pt,topsep=2pt]
  \item Try It 2a
  \item Try It 2b
  \item Practice \#14
  \item Practice \#15
  \item Practice \#16
\end{enumerate}

\bankitem{Try It 2a}{%
Describe the behavior of the graph of the function at each of its zeros: \(f(x) = x(x+4)(x-1)^4.\)
}{}

\bankitem{Try It 2b}{%
Describe the behavior of the graph of the function at each of its zeros: \(f(x) = (x^2 + 9)(x - 1)^5(x + 2)^2.\)
}{}

\bankitem{Practice \#14}{%
Find the zeros of the function, and describe the behavior of the graph at each zero: \(f(x) = x^3 - 8x^2 + 16x.\)
}{}

\bankitem{Practice \#15}{%
Find the zeros of the function, and describe the behavior of the graph at each zero: \(g(x) = x^3 - x^2 - 25x + 25.\)
}{}

\bankitem{Practice \#16}{%
Find the zeros of the function, and describe the behavior of the graph at each zero: \(f(x) = 9x^4 - 40x^2 + 16.\)
}{}

\sectionbanner{SHARE / SUMMARY (5 min)}

\begin{sentenceframebox}
\textbf{Sentence frames}

Pair share: ``For f(x) = \_\_\_, the zero x = \_\_\_ has multiplicity \_\_\_, so the graph \_\_\_ the x-axis at that zero.''

Whole class: one pair at random reads their sentence. Class signals thumbs up/sideways.
\end{sentenceframebox}

\begin{callout}{\IconSearch\ BRIDGE TO TOMORROW}{calloutpeach}
If a cubic has zeros 2, \(1 + \sqrt{3}\), and \(1 - \sqrt{3}\), what is the nature of its coefficients --- rational, irrational, or complex? (Tuesday-to-Wednesday bridge; revisit at start of Period 3.)
\end{callout}

\sectionbanner{EXIT}

\begin{summaryexitbox}{EXIT TICKET --- Summary Recap}
Today I learned that \dotfill

The rule I want to remember is \dotfill

One thing I'm still unsure about is \dotfill
\end{summaryexitbox}

\clearpage

\sectionbanner{OPTIONAL REINFORCEMENT (not collected)}
\textit{If you finish early or want extra practice before Wednesday. Answers on Desmos or ask for a check.}

\bankitem{Optional \#1\ \ (Savvas Practice)}{%
At what points do the graphs of \(f(x) = x^3 - 2x^2 - 16x + 20\) and \(g(x) = -12\) intersect?
}{}

\bankitem{Optional \#2\ \ (Savvas Practice)}{%
Sketch the graph of the function by finding the zeros: \(f(x) = 3x^3 - 9x^2 - 12x.\)
}{}

\end{document}
